% slides.tex -- lab meeting deck (about 20 slides + backup, ~30 min)
% Build:  python make_slide_assets.py   (from the scripts folder)
%         cd slides && pdflatex slides.tex && pdflatex slides.tex
% Numbers in the text come from tables/numbers.tex via \res{...}; '??' means a run folder was missing.
\documentclass[aspectratio=169,11pt]{beamer}
\usetheme{Madrid}
\setbeamertemplate{navigation symbols}{}
\usepackage[T1]{fontenc}
\usepackage{lmodern}
\usepackage{amsmath,amssymb,booktabs,graphicx}

\makeatletter
\newcommand{\res}[1]{\@ifundefined{res@#1}{\textbf{??}}{\@nameuse{res@#1}}}
\makeatother
\newcommand{\safeinput}[1]{\IfFileExists{#1}{\input{#1}}{\fbox{\footnotesize missing: \texttt{\detokenize{#1}}}}}
\newcommand{\safegraphic}[2]{\IfFileExists{#2}{\includegraphics[#1]{#2}}{\fbox{\parbox{0.5\linewidth}{\centering\footnotesize missing: \texttt{\detokenize{#2}}}}}}
\IfFileExists{tables/numbers.tex}{\makeatletter
\@namedef{res@A-aucb-ALL}{0.625}
\@namedef{res@A-aucb-gaussian}{0.669}
\@namedef{res@A-aucb-hotspot}{0.543}
\@namedef{res@A-aucb-linear}{0.980}
\@namedef{res@A-aucb-periodic}{0.527}
\@namedef{res@A-aucbeta}{0.625}
\@namedef{res@A-aucp-ALL}{0.505}
\@namedef{res@A-aucp-gaussian}{0.511}
\@namedef{res@A-aucp-hotspot}{0.477}
\@namedef{res@A-aucp-linear}{0.876}
\@namedef{res@A-aucp-periodic}{0.391}
\@namedef{res@A-aucpip}{0.505}
\@namedef{res@A-called}{673}
\@namedef{res@A-f1}{0.141}
\@namedef{res@A-fn}{77}
\@namedef{res@A-fp}{594}
\@namedef{res@A-ngenes}{1,500}
\@namedef{res@A-ntrue}{132}
\@namedef{res@A-nunits}{1,316}
\@namedef{res@A-prec}{0.085}
\@namedef{res@A-rate-adip}{42\%}
\@namedef{res@A-rate-ctsv}{73\%}
\@namedef{res@A-rate-marker}{47\%}
\@namedef{res@A-rate-null}{44\%}
\@namedef{res@A-rate-svgaussian}{33\%}
\@namedef{res@A-rate-svhotspot}{41\%}
\@namedef{res@A-rate-svlinear}{94\%}
\@namedef{res@A-rate-svperiodic}{30\%}
\@namedef{res@A-rec}{0.417}
\@namedef{res@A-runtime}{23.0}
\@namedef{res@A-topk}{0.250}
\@namedef{res@A-tp}{55}
\@namedef{res@B-aucb-ALL}{0.569}
\@namedef{res@B-aucb-gaussian}{0.629}
\@namedef{res@B-aucb-hotspot}{0.553}
\@namedef{res@B-aucb-linear}{0.930}
\@namedef{res@B-aucb-periodic}{0.410}
\@namedef{res@B-aucbeta}{0.569}
\@namedef{res@B-aucp-ALL}{0.507}
\@namedef{res@B-aucp-gaussian}{0.532}
\@namedef{res@B-aucp-hotspot}{0.492}
\@namedef{res@B-aucp-linear}{0.806}
\@namedef{res@B-aucp-periodic}{0.393}
\@namedef{res@B-aucpip}{0.507}
\@namedef{res@B-called}{859}
\@namedef{res@B-f1}{0.149}
\@namedef{res@B-fn}{60}
\@namedef{res@B-fp}{761}
\@namedef{res@B-ngenes}{1,500}
\@namedef{res@B-ntrue}{132}
\@namedef{res@B-nunits}{1,316}
\@namedef{res@B-prec}{0.086}
\@namedef{res@B-rate-adip}{56\%}
\@namedef{res@B-rate-ctsv}{79\%}
\@namedef{res@B-rate-marker}{58\%}
\@namedef{res@B-rate-null}{57\%}
\@namedef{res@B-rate-svgaussian}{58\%}
\@namedef{res@B-rate-svhotspot}{53\%}
\@namedef{res@B-rate-svlinear}{94\%}
\@namedef{res@B-rate-svperiodic}{39\%}
\@namedef{res@B-rec}{0.545}
\@namedef{res@B-runtime}{2.1}
\@namedef{res@B-topk}{0.167}
\@namedef{res@B-tp}{72}
\@namedef{res@C-aucb-ALL}{0.760}
\@namedef{res@C-aucb-gaussian}{0.874}
\@namedef{res@C-aucb-hotspot}{0.659}
\@namedef{res@C-aucb-linear}{0.998}
\@namedef{res@C-aucb-periodic}{0.662}
\@namedef{res@C-aucbeta}{0.760}
\@namedef{res@C-aucp-ALL}{0.734}
\@namedef{res@C-aucp-gaussian}{0.773}
\@namedef{res@C-aucp-hotspot}{0.643}
\@namedef{res@C-aucp-linear}{0.900}
\@namedef{res@C-aucp-periodic}{0.712}
\@namedef{res@C-aucpip}{0.734}
\@namedef{res@C-called}{785}
\@namedef{res@C-f1}{0.246}
\@namedef{res@C-fn}{22}
\@namedef{res@C-fp}{651}
\@namedef{res@C-ngenes}{1,500}
\@namedef{res@C-ntrue}{132}
\@namedef{res@C-nunits}{1,316}
\@namedef{res@C-prec}{0.145}
\@namedef{res@C-rate-adip}{55\%}
\@namedef{res@C-rate-ctsv}{73\%}
\@namedef{res@C-rate-marker}{48\%}
\@namedef{res@C-rate-null}{44\%}
\@namedef{res@C-rate-svgaussian}{89\%}
\@namedef{res@C-rate-svhotspot}{65\%}
\@namedef{res@C-rate-svlinear}{100\%}
\@namedef{res@C-rate-svperiodic}{87\%}
\@namedef{res@C-rec}{0.833}
\@namedef{res@C-runtime}{1.9}
\@namedef{res@C-topk}{0.364}
\@namedef{res@C-tp}{110}
\@namedef{res@D-aucb-ALL}{0.837}
\@namedef{res@D-aucb-gaussian}{0.915}
\@namedef{res@D-aucb-hotspot}{0.681}
\@namedef{res@D-aucb-linear}{0.997}
\@namedef{res@D-aucb-periodic}{0.836}
\@namedef{res@D-aucbeta}{0.837}
\@namedef{res@D-aucp-ALL}{0.705}
\@namedef{res@D-aucp-gaussian}{0.699}
\@namedef{res@D-aucp-hotspot}{0.562}
\@namedef{res@D-aucp-linear}{0.865}
\@namedef{res@D-aucp-periodic}{0.761}
\@namedef{res@D-aucpip}{0.705}
\@namedef{res@D-called}{525}
\@namedef{res@D-f1}{0.257}
\@namedef{res@D-fn}{49}
\@namedef{res@D-fp}{431}
\@namedef{res@D-ngenes}{1,500}
\@namedef{res@D-ntrue}{132}
\@namedef{res@D-nunits}{1,316}
\@namedef{res@D-prec}{0.161}
\@namedef{res@D-rate-adip}{36\%}
\@namedef{res@D-rate-ctsv}{33\%}
\@namedef{res@D-rate-marker}{30\%}
\@namedef{res@D-rate-null}{31\%}
\@namedef{res@D-rate-svgaussian}{61\%}
\@namedef{res@D-rate-svhotspot}{38\%}
\@namedef{res@D-rate-svlinear}{88\%}
\@namedef{res@D-rate-svperiodic}{74\%}
\@namedef{res@D-rec}{0.629}
\@namedef{res@D-runtime}{9.4}
\@namedef{res@D-topk}{0.439}
\@namedef{res@D-tp}{83}
\@namedef{res@E-aucb-ALL}{0.809}
\@namedef{res@E-aucb-gaussian}{0.915}
\@namedef{res@E-aucb-hotspot}{0.670}
\@namedef{res@E-aucb-linear}{0.997}
\@namedef{res@E-aucb-periodic}{0.765}
\@namedef{res@E-aucbeta}{0.809}
\@namedef{res@E-aucp-ALL}{0.732}
\@namedef{res@E-aucp-gaussian}{0.804}
\@namedef{res@E-aucp-hotspot}{0.627}
\@namedef{res@E-aucp-linear}{0.786}
\@namedef{res@E-aucp-periodic}{0.735}
\@namedef{res@E-aucpip}{0.732}
\@namedef{res@E-called}{300}
\@namedef{res@E-f1}{0.317}
\@namedef{res@E-fn}{65}
\@namedef{res@E-fp}{224}
\@namedef{res@E-ngenes}{1,500}
\@namedef{res@E-ntrue}{132}
\@namedef{res@E-nunits}{1,316}
\@namedef{res@E-prec}{0.230}
\@namedef{res@E-rate-adip}{18\%}
\@namedef{res@E-rate-ctsv}{27\%}
\@namedef{res@E-rate-marker}{16\%}
\@namedef{res@E-rate-null}{16\%}
\@namedef{res@E-rate-svgaussian}{58\%}
\@namedef{res@E-rate-svhotspot}{32\%}
\@namedef{res@E-rate-svlinear}{69\%}
\@namedef{res@E-rate-svperiodic}{52\%}
\@namedef{res@E-rec}{0.508}
\@namedef{res@E-runtime}{9.7}
\@namedef{res@E-topk}{0.424}
\@namedef{res@E-tp}{67}
\@namedef{res@F-aucbeta}{0.999}
\@namedef{res@F-aucpip}{0.996}
\@namedef{res@F-called}{590}
\@namedef{res@F-f1}{0.917}
\@namedef{res@F-fn}{0}
\@namedef{res@F-fp}{90}
\@namedef{res@F-fphub}{0.191}
\@namedef{res@F-lam}{0.0826}
\@namedef{res@F-lam23}{0.190}
\@namedef{res@F-ngenes}{5,000}
\@namedef{res@F-nnonhub}{4,410}
\@namedef{res@F-nnullhub}{90}
\@namedef{res@F-nonhub}{17}
\@namedef{res@F-ntrue}{500}
\@namedef{res@F-nullhub}{73}
\@namedef{res@F-nunits}{1,024}
\@namedef{res@F-prec}{0.847}
\@namedef{res@F-rec}{1.000}
\@namedef{res@F-runtime}{16.3}
\@namedef{res@F-topk}{0.978}
\@namedef{res@F-tp}{500}
\@namedef{res@G-aucbeta}{1.000}
\@namedef{res@G-aucpip}{0.976}
\@namedef{res@G-called}{953}
\@namedef{res@G-f1}{0.688}
\@namedef{res@G-fn}{0}
\@namedef{res@G-fp}{453}
\@namedef{res@G-fphub}{0.055}
\@namedef{res@G-lam}{0.0367}
\@namedef{res@G-lam23}{0.111}
\@namedef{res@G-ngenes}{5,000}
\@namedef{res@G-nnonhub}{4,410}
\@namedef{res@G-nnullhub}{90}
\@namedef{res@G-nonhub}{434}
\@namedef{res@G-ntrue}{500}
\@namedef{res@G-nullhub}{19}
\@namedef{res@G-nunits}{1,024}
\@namedef{res@G-prec}{0.525}
\@namedef{res@G-rec}{1.000}
\@namedef{res@G-runtime}{16.1}
\@namedef{res@G-topk}{0.998}
\@namedef{res@G-tp}{500}
\@namedef{res@fig-nbins}{101,000}
\@namedef{res@fig-ngenes}{18,000}
\@namedef{res@fit-nbins}{20,000}
\@namedef{res@fit-ngenes}{4,000}
\@namedef{res@ours-cpg}{0.37}
\@namedef{res@paper-cpg}{20.5}
\@namedef{res@paper-f1}{0.996}
\@namedef{res@paper-nbins}{1,024}
\@namedef{res@paper-ngenes}{5,000}
\@namedef{res@paper-nsv}{500}
\@namedef{res@paper-prec}{1.000}
\@namedef{res@paper-rec}{0.992}
\@namedef{res@paper-scen}{linear pattern, star network, dropout 0.1}
\makeatother
}{}

\title[NABR for spatially variable genes]{Network-assisted Bayesian detection of spatially variable genes}
\subtitle{Implementation, synthetic benchmark, and partial replication of Wu et al.\ (2025, \emph{AOAS})}
\author{Alireza}   % <-- edit
\institute{Lab meeting}
\date{\today}

\begin{document}

\begin{frame}
  \titlepage
\end{frame}

\begin{frame}{Outline}
  \begin{enumerate}
    \item Problem and method: network-assisted Bayesian regularization (NABR)
    \item Our implementation
    \item Synthetic benchmark resembling our data
    \item Replication on the paper's own simulation design
    \item Diagnosis, take-home messages, next steps
  \end{enumerate}
\end{frame}

% =====================================================================
\section{Problem and method}

\begin{frame}{The problem}
  \begin{itemize}
    \item Binned spatial transcriptomics (16\,$\mu$m bins); the \textbf{bin type is known} for every bin.
    \item Goal: \textbf{global spatially variable (SV) genes} --- expression that varies across the tissue
          beyond what bin type explains.
    \item Pitfall: compartments are spatially organized $\Rightarrow$ compartment markers look ``spatial''.
          Bin type must be adjusted for.
    \item Counts are sparse, over-dispersed and zero-inflated; genes act in networks.
  \end{itemize}
\end{frame}

\begin{frame}{Method: Wu, Li, Ma \& Wu (2025)}
  \small Joint identification of spatially variable genes via a network-assisted Bayesian regularization
  approach. \emph{Ann.\ Appl.\ Stat.} 19(4):2705--2723.
  \normalsize
  \begin{enumerate}
    \item \textbf{ZINB counts + cell-type composition covariates} $\Rightarrow$ corrects composition confounding.
    \item \textbf{Mean-based spatial effects} with fixed kernels $\Rightarrow$ no $n\times n$ covariance, scalable.
    \item \textbf{Network-assisted thresholded Gaussian prior} $\Rightarrow$ \emph{joint} (not marginal)
          selection of SV genes.
  \end{enumerate}
  \vspace{2mm}
  In the paper's 12 simulation scenarios: F1 between 0.91 and 1.00, the highest among SPARK, SPARK-X,
  SpatialDE, MERINGUE, CTSV and HEARTSVG.
\end{frame}

\begin{frame}{Model}
  \begin{align*}
    Y_{ij} &\sim (1-r_{ij})\,\mathrm{NB}(c_i\mu_{ij},\,\phi) + r_{ij}\,\delta_0,
           \qquad r_{ij}\sim\mathrm{Bern}(\pi_j)\\[2mm]
    \log\mu_{ij} &= \mu_{0j} + \beta^{(1)}_j\,\mathcal K(x_{i1}) + \beta^{(2)}_j\,\mathcal K(x_{i2})
           + \textstyle\sum_{k=1}^{K}\alpha_k\,w_{ik}
  \end{align*}
  \begin{itemize}
    \item $w_{ik}$: cell-type composition of unit $i$ (for us: fraction of bins of type $k$).
    \item $\mathcal K$: linear, exponential or periodic function of standardized coordinates.
    \item Gene $j$ is SV $\iff \beta^{(1)}_j\neq0$ or $\beta^{(2)}_j\neq0$.
  \end{itemize}
\end{frame}

\begin{frame}{Prior: sparsity + network}
  \begin{align*}
    \boldsymbol\beta &= \boldsymbol\gamma\circ\mathbf 1\{|\boldsymbol\gamma|>\lambda\boldsymbol\rho\},
    &\boldsymbol\gamma &\sim N\big(\mathbf 0,\ \sigma^2_\gamma(\mathbf L+\varepsilon\mathbf I)^{-1}\big),\\
    \mathbf L &= \big(\mathrm{sgn}\tilde{\boldsymbol\beta}\,\mathrm{sgn}\tilde{\boldsymbol\beta}^{\top}\big)\circ
       \big(\mathbf I-\mathbf D^{-1/2}\mathbf A\,\mathbf D^{-1/2}\big),
    &\rho_j &= \Big(\max_{l\in V_{s(j)}}|\gamma_l|\Big)^{-1/2},
  \end{align*}
  \centerline{$\lambda\sim U\big(0,\ q_{0.9}(|\tilde{\boldsymbol\beta}|)\big)$, \quad
  $\sigma^2_\gamma\sim IG(3.5,\,0.5)$.}
  \begin{itemize}
    \item Hard threshold $\Rightarrow$ exact zeros; $\lambda$ acts as the minimal detectable effect.
    \item Graph Laplacian $\Rightarrow$ connected genes get similar effect magnitudes.
    \item Group-wise weights $\rho$ lower the threshold in sub-networks that contain a strong gene.
  \end{itemize}
\end{frame}

\begin{frame}{Inference and gene calling}
  \begin{itemize}
    \item MCMC: Gibbs ($r_{ij}$, $\sigma^2_\gamma$); Metropolis--Hastings ($\mu_{0j}$, $\alpha$, $\phi$, $\lambda$);
          preconditioned Crank--Nicolson Langevin ($\boldsymbol\gamma$); deterministic update of $\rho$.
    \item One run per kernel (linear, 2 exponential, 2 cosine); combined with posterior model probabilities.
    \item PIP$_j$ = posterior probability that gene $j$ is selected; Bayesian FDR at 0.05.
    \item 5 chains; genes selected in $\geq 80\%$ of chains are reported.
  \end{itemize}
\end{frame}

% =====================================================================
\section{Implementation}

\begin{frame}{Our implementation}
  \begin{itemize}
    \item Python / PyTorch, all MCMC steps (a)--(h) of the paper; CPU or GPU.
    \item About 1\,s per iteration for 1{,}024 spots $\times$ 5{,}000 genes on a 4-core laptop CPU.
    \item Pipeline: simulate $\rightarrow$ fit $\rightarrow$ evaluate $\rightarrow$ plot.
    \item Documented deviations (each switchable):
      \begin{itemize}
        \item gene-specific bin-type effects $\alpha_{jk}$ for our data (the paper's shared $\alpha_k$ is used in the replication);
        \item exact pCN-Langevin acceptance ratio (target and proposal densities);
        \item kernel columns standardized; $\pi_j$ sampled; proposal tuning only during burn-in.
      \end{itemize}
  \end{itemize}
\end{frame}

% =====================================================================
\section{Our synthetic benchmark}

\begin{frame}{Synthetic tissue resembling our data}
  \begin{columns}[T]
    \begin{column}{0.58\textwidth}
      \safegraphic{width=\linewidth,height=0.75\textheight,keepaspectratio}{figs/tissue.png}
    \end{column}
    \begin{column}{0.40\textwidth}
      \small
      \begin{itemize}
        \item \res{fig-nbins} bins of 16\,$\mu$m, \res{fig-ngenes} genes (shown).
        \item 6 compartments; adipose has low RNA and $\sim$70\% of genes switched off.
        \item ZINB counts with gene-specific dispersion and dropout.
        \item Gene network: star / scale-free modules, singletons, random edges.
      \end{itemize}
    \end{column}
  \end{columns}
\end{frame}

\begin{frame}{Sequencing depth and detected genes by compartment}
  \centering
  \safegraphic{width=\linewidth,height=0.8\textheight,keepaspectratio}{figs/qc.png}
\end{frame}

\begin{frame}{Ground truth}
  \centering
  \safegraphic{width=\linewidth,height=0.62\textheight,keepaspectratio}{figs/truth.png}
  \small
  \begin{itemize}
    \item SV genes: linear, Gaussian, periodic and hotspot patterns, mostly inside network modules.
    \item Decoys: compartment markers and adipose-off genes (spatially structured but \emph{null});
          cell-type-specific SV genes (reported separately).
  \end{itemize}
\end{frame}

\begin{frame}{Results on our synthetic data}
  \centering
  \resizebox{\linewidth}{!}{\safeinput{tables/runs_ours.tex}}
  \vspace{2mm}
  \small
  \begin{itemize}
    \item Fits used a \res{fit-nbins}-bin, \res{fit-ngenes}-gene version; $4\times4$-bin pseudo-spots
          (\res{B-nunits} units), \res{B-ngenes} highly variable genes; short MCMC (400 iterations).
    \item BFDR calls are not usable here: precision between \res{A-prec} and \res{E-prec}.
  \end{itemize}
\end{frame}

\begin{frame}{Diagnosis 1: bin-type confounding is not the problem}
  \centering
  \safegraphic{width=\linewidth,height=0.62\textheight,keepaspectratio}{figs/callrate.pdf}
  \small
  \begin{itemize}
    \item Null genes are called at similar rates whatever their type (run B: markers \res{B-rate-marker},
          adipose-off \res{B-rate-adip}, other nulls \res{B-rate-null}).
    \item Over-calling is generic, not leakage of compartment structure.
  \end{itemize}
\end{frame}

\begin{frame}{Diagnosis 2: ranking is informative, calling is not}
  \centering
  \safegraphic{width=\linewidth,height=0.6\textheight,keepaspectratio}{figs/auc_pattern.pdf}
  \small
  \begin{itemize}
    \item Ranking by posterior $|\hat\beta|$ (run D): linear \res{D-aucb-linear}, Gaussian \res{D-aucb-gaussian},
          periodic \res{D-aucb-periodic}, hotspot \res{D-aucb-hotspot} (hotspots lie outside the kernel family).
    \item With $\rho$ on (run A) the PIP ranking degrades.
  \end{itemize}
\end{frame}

\begin{frame}{Precision of the top-ranked genes}
  \centering
  \safegraphic{width=\linewidth,height=0.78\textheight,keepaspectratio}{figs/topk.pdf}
\end{frame}

% =====================================================================
\section{Replication}

\begin{frame}{Replicating the paper's simulation design}
  \centering
  \footnotesize
  \begin{tabular}{p{3.0cm}p{5.2cm}p{5.6cm}}
    \toprule
    & Paper design (replicated) & Our design \\
    \midrule
    Units analysed & \res{paper-nbins} spots ($32\times32$) & \res{B-nunits} pseudo-spots ($4\times4$ bins) \\
    Genes analysed & \res{paper-ngenes} (all) & \res{B-ngenes} HVGs of \res{fit-ngenes} \\
    Counts/gene/unit (median) & \res{paper-cpg} & \res{ours-cpg} \\
    Bin-type information & Dirichlet compositions, 3 regions & hard labels, 6 irregular compartments \\
    Cell-type effect & shared $\alpha_k\sim N(0,3.5^2)$ & gene-specific: markers, genes off in adipose \\
    Spatial patterns & linear (inside the kernel family) & linear, Gaussian, periodic, hotspot \\
    Network & 100 sub-networks $\times$ 50 genes; all SV in 10 & modules + singletons + random edges;
              20\% nulls inside SV modules \\
    Decoys & none & markers, adipose-off, cell-type-specific SV \\
    \bottomrule
  \end{tabular}
\end{frame}

\begin{frame}{Replication data}
  \centering
  \safegraphic{width=\linewidth,height=0.45\textheight,keepaspectratio}{figs/paper_design.pdf}
  \small
  \begin{itemize}
    \item Scenario: \res{paper-scen}; \res{paper-nbins} spots $\times$ \res{paper-ngenes} genes;
          \res{paper-nsv} SV genes in 10 star sub-networks.
    \item Region layout, coordinate scaling and scale-free generator are not specified in the paper:
          our choices.
  \end{itemize}
\end{frame}

\begin{frame}{Replication results}
  \centering
  \resizebox{\linewidth}{!}{\safeinput{tables/runs_paper.tex}}
  \vspace{2mm}
  \small
  \begin{itemize}
    \item All \res{F-ntrue} SV genes recovered (recall \res{F-rec}); ranking by $|\hat\beta|$:
          AUROC \res{F-aucbeta} ($\rho$ on), \res{G-aucbeta} ($\rho$ off).
    \item Precision \res{F-prec} vs.\ \res{paper-prec} in the paper. \textbf{Partial} replication: one replicate,
          linear kernel, one chain, 1{,}000 iterations (paper: 50 replicates, 5 kernels $\times$ 5 chains
          $\times$ 2{,}000 iterations).
  \end{itemize}
\end{frame}

\begin{frame}{Where are the false positives? Ablation of $\rho$}
  \centering
  \safegraphic{width=0.85\linewidth,height=0.55\textheight,keepaspectratio}{figs/hub_ablation.pdf}
  \small
  \begin{itemize}
    \item $\rho$ on: \res{F-nullhub}/\res{F-nnullhub} null hubs and \res{F-nonhub}/\res{F-nnonhub} other nulls called.
    \item $\rho$ off: \res{G-nullhub}/\res{G-nnullhub} null hubs and \res{G-nonhub}/\res{G-nnonhub} other nulls called.
    \item $\Rightarrow$ $\rho$ is needed to silence noise in ordinary genes, but promotes the hubs of null star networks.
  \end{itemize}
\end{frame}

\begin{frame}{A candidate mechanism (hypothesis, not yet tested)}
  \begin{columns}[T]
    \begin{column}{0.55\textwidth}
      \safegraphic{width=\linewidth,height=0.7\textheight,keepaspectratio}{figs/mechanism.pdf}
    \end{column}
    \begin{column}{0.43\textwidth}
      \small
      If the hub $h$ has the largest $|\gamma|$ in its sub-network, $\rho_h=|\gamma_h|^{-1/2}$, so
      \[ h \text{ selected} \iff |\gamma_h|>\lambda\rho_h \iff |\gamma_h|>\lambda^{2/3}. \]
      Observed: $\lambda=\res{F-lam}$, $\lambda^{2/3}=\res{F-lam23}$;
      called null hubs: median $|\hat\beta|=\res{F-fphub}$.
      \begin{itemize}
        \item $\rho$ is recomputed without a likelihood check.
        \item The Laplacian prior is nearly flat along $\mathbf D^{1/2}\mathbf 1$, where a star hub
              carries half the weight.
      \end{itemize}
    \end{column}
  \end{columns}
\end{frame}

\begin{frame}{Network view (replication, $\rho$ on)}
  \centering
  \safegraphic{width=\linewidth,height=0.8\textheight,keepaspectratio}{figs/network_paper.png}
\end{frame}

% =====================================================================
\section{Conclusions}

\begin{frame}{Take-home messages}
  \begin{enumerate}
    \item On the paper's design the implementation recovers \res{F-tp}/\res{F-ntrue} SV genes, with
          AUROC$_{|\hat\beta|}$ = \res{F-aucbeta}.
    \item BFDR calls are not yet calibrated (precision \res{F-prec} vs.\ \res{paper-prec});
          errors concentrate in null star hubs $\Rightarrow$ a specific, testable mechanism.
    \item $\rho$ is a trade-off: it silences noise in ordinary genes but promotes null hubs.
    \item On sparse data (median \res{ours-cpg} vs.\ \res{paper-cpg} counts per gene per unit) performance drops
          sharply; bin-type confounding itself is handled.
    \item Promising for our data, but not yet validated at Visium-HD-like sparsity.
  \end{enumerate}
\end{frame}

\begin{frame}{Recommendations for our data now}
  \begin{itemize}
    \item Aggregate bins (32--64\,$\mu$m) to raise counts per analysis unit.
    \item Rank genes by posterior effect size; treat BFDR calls as provisional.
    \item Calibrate calls with a permutation null (shuffle coordinates within bin type).
    \item Benchmark against nnSVG and SPARK-X (bin type as covariate) on the same simulation.
  \end{itemize}
\end{frame}

\begin{frame}{Next steps}
  \begin{enumerate}
    \item Full replication: 12 scenarios $\times$ 5 kernels $\times$ 5 chains (HPC CPU array).
    \item Test the hub mechanism: likelihood-checked $\rho$ update, degree-aware weights.
    \item Bridge the two designs one factor at a time: counts, gene-specific bin-type effects,
          kernel mismatch, network noise, hard labels.
    \item Apply to our real data.
  \end{enumerate}
\end{frame}

\begin{frame}
  \centering \Large Thank you --- questions?
\end{frame}

% =====================================================================
\appendix

\begin{frame}{Backup: MCMC traces (replication)}
  \centering
  \safegraphic{width=\linewidth,height=0.5\textheight,keepaspectratio}{figs/traces_paper.pdf}
  \small
  \begin{itemize}
    \item $\phi$ converges slowly (shared $\alpha_k$ start at zero) and reaches the true value $\approx 10$
          only after about 650 iterations.
  \end{itemize}
\end{frame}

\begin{frame}{Backup: replication summary figure ($\rho$ on)}
  \centering
  \safegraphic{width=\linewidth,height=0.8\textheight,keepaspectratio}{figs/results_paper.png}
\end{frame}

\begin{frame}{Backup: gene maps on our data (run D)}
  \centering
  \safegraphic{width=\linewidth,height=0.8\textheight,keepaspectratio}{figs/maps_ours.png}
\end{frame}

\begin{frame}[fragile]{Backup: reproducibility}
\footnotesize
\begin{verbatim}
python main.py --stage all --quick --sim-bins 20000 --sim-genes 4000      # run A
python debug_norho.py; python debug_norho_all.py; python debug_floor.py   # runs C, D, E
python sim_paper.py --outdir sim_paper
python run_replication.py --sim-dir sim_paper --fit-dir fit_paper --kernels linear \
       --n-chains 1 --n-iter 1000 --burn-in 500 --thin 5                  # run F
python run_replication.py ... --fit-dir fit_paper_norho --no-rho          # run G
python make_slide_assets.py; cd slides; pdflatex slides.tex
\end{verbatim}
\end{frame}

\end{document}
